\chapter{Matrix-free operators}
\label{ch:matrixfree}
\impl{src/physics.jl, src/two\_phase\_action.jl}

\section{The action, not the matrix}

Conjugate gradients never needs $\Kmat$ --- only its action $\Kmat\vect{v}$ on
a sequence of vectors.  A matrix-free operator supplies that action by
assembling directly:
\begin{equation}
  \left(\Kmat\vect{v}\right)_a
  = \int_{\Omega_0} \Grad N_a \cdot \fourth{A} \cdot \Grad N_b \, \mathrm{d}V \ v_b
  = \int_{\Omega_0} \Grad N_a \dbldot \delta\Piola(\Grad \vect{v}) \, \mathrm{d}V,
  \label{eq:action}
\end{equation}
where $\delta\Piola$ is the directional derivative of the stress in the
direction $\Grad\vect{v}$ --- equation~\eqref{eq:nh-action} for neo-Hookean.
The right-hand form is the one that matters: the tangent's action on
\emph{one} direction never requires the $81$-component $\fourth{A}$ to be
formed, so the closed-form derivative costs a few $3\times3$ products per
quadrature point.

Three things follow.  Memory drops from $O(\text{nnz})$ to $O(n)$ --- measured
at $5.83$~GB down to $0.245$~GB on a $530$k-DOF problem, which is the
difference between fitting on a GPU and not.  There is no assembly phase and no
sparse pattern to build.  And the operator is exactly as current as the
displacement it is evaluated at, so no staleness question arises.

The cost is that each application redoes the constitutive work.  Whether that
is a good trade depends on the arithmetic intensity of the element kernel
against the bandwidth of storing and re-reading a sparse matrix; on GPU it is
decisively favourable.

\section{Operator fusion}

The Newmark system needs $\Keff\vect{v} = \Kmat\vect{v} + c_M\Mmat\vect{v}$.
Computed as two separate assemblies, the element loop runs twice and each pass
gathers the same nodal data.  Fusing them into a single kernel that accumulates
both contributions per quadrature point halves the gather traffic, and the
gather is what the operator is bound by.

The same reasoning produces a diagonal-only kernel for the Jacobi
preconditioner: $\diag(\Keff)$ is assembled directly rather than extracted from
an operator that is never formed.

\section{Where the time goes}

Profiling the fused Newmark action on an A100 at $530$k DOF attributes $67\%$
of device time to the stiffness matrix--vector product, with the device busy
$72\%$ of the step.  The step, to a good approximation, \emph{is} the matvec.

That single fact organizes the optimization work.  It is why atomic contention
is not the bottleneck --- ablating the atomics changes nothing measurable ---
and why the loss ordering across architectures tracks effective cache size
rather than arithmetic throughput.  It is also why the forcing terms of
Chapter~\ref{ch:inexact} pay so much better on the GPU paths than on assembled
CPU runs: on this path, removing a conjugate-gradient iteration removes almost
a whole iteration's worth of wall time.

\section{Two-phase assembly}

The default kernel assigns one thread per element: gather the element's nodal
values, loop the quadrature points serially, scatter the results with atomic
adds.  The working set held across that quadrature loop is what limits
occupancy.

The two-phase form splits it.  Phase one assigns one thread per
(element, quadrature point) pair, writing each contribution to a disjoint slot
of a staging array --- eight times the threads, an eighth of the per-thread
state, and no atomics.  Phase two assigns one thread per node and sums the
staging slots incident on it through an inverse adjacency built once on the
host, again without atomics.

The transformation is sound and it is implemented, but measurement did not
support it: $0.66$ to $0.80$ times the throughput of the one-thread-per-element
kernel on all three architectures tested.  The extra staging traffic exceeds
what the occupancy gain returns.  It is documented here because the reasoning
is correct and the conclusion is specific to the current data layout --- an
element-blocked field layout could change it.

\section{Mixed precision}
\label{sec:mixed-precision}

The V-cycle of Chapter~\ref{ch:amg} applies the matrix-free action five times
per call at $\nu = 2$, which on a consumer GPU --- where double precision runs
at $1/32$ rate --- dominates the iteration.  Carina evaluates the fine-level
\emph{smoother} action in single precision while the operator conjugate
gradients solves against stays in double.

This is legitimate, and the argument is the one from
Remark~\ref{ch:constitutive}.  A preconditioner may be any symmetric positive
definite approximation; reducing its precision changes which approximation it
is, not what the iteration converges to.  What would \emph{not} be legitimate
is reducing the precision of the operator itself, which would change the fixed
point.  The reduced-precision V-cycle also remains a fixed linear operator
across applications, so ordinary conjugate gradients --- rather than a flexible
variant --- stays valid.

The gate is a property of the model, not a user request: Carina enables the
reduced-precision path only for models that demonstrably narrow to single
precision, and verifies this rather than assuming it.  A model whose action
silently falls back to double would otherwise present as a preconditioner that
was never actually accelerated.
